% fig_vl09_unschaerfe.tex — Woher die Schranke kommt: die Positivität von
% M = C^dag C ist eine Parabel in lambda, die die Nulllinie nicht
% unterschreiten darf, und genau ihre Scheitelhöhe ist Kennards
% Ungleichung.
%
% Follows the conventions of fig_vl11_kasten.tex.
%
% WARUM NICHT DIE GAUSS-PAKETE.  Die naheliegende Tafel — schmales
% Ortspaket, breites Impulspaket — steht bereits in vl08
% (fig_vl08_fourier): dort sind drei Pakete in beiden Darstellungen
% nebeneinandergestellt und mit sigma_x sigma_p = hbar/2 beschriftet.
% Sie hier zu wiederholen hieße, die vorige Vorlesung noch einmal zu
% zeichnen.  Neu in vl09 ist der *Beweis*: der Kunstgriff mit M, das
% Minimieren über lambda, und was daraus für zulässige Paare
% (Delta x, Delta p) folgt.  Genau das ist hier gezeichnet — Panel (a)
% der Mechanismus, Panel (b) sein Ergebnis.
%
% GEOMETRIE AUS DER FORMEL.  Alles folgt aus der Ungleichung der
% Vorlesung, mit A = p - <p> und B = x - <x>:
%
%       tr(rho M) = tr(rho A^2) + lambda^2 tr(rho B^2)
%                                       + lambda tr(rho i[A,B]) >= 0 .
%
% Hier ist tr(rho A^2) = (Delta p)^2, tr(rho B^2) = (Delta x)^2 und,
% wegen [x,p] = i hbar, tr(rho i[A,B]) = i[p,x] = hbar.  Also
%
%       f(lambda) = (Delta p)^2 + lambda^2 (Delta x)^2 + lambda hbar .
%
% In den dimensionslosen Variablen
%
%       mu    = lambda (Delta x)^2 / hbar ,
%       g(mu) = f (Delta x)^2 / hbar^2 ,
%       kappa = Delta x Delta p / hbar
%
% lautet das
%
%       g(mu) = mu^2 + mu + kappa^2 ,
%
% eine Parabelschar, die NUR vom Produkt kappa abhängt — die einzige
% dimensionslose Größe des Problems.  Scheitel bei mu = -1/2, das heißt
% lambda_* = -hbar / 2 (Delta x)^2, genau der Wert, den die Vorlesung als
% Übung ausrechnen lässt; Scheitelhöhe kappa^2 - 1/4.  Positivität heißt
% also kappa >= 1/2, und das ist Delta x Delta p >= hbar/2.  Die drei
% Kurven sind demnach nicht gewählt, sondern abgelesen: kappa = 1
% (Ungleichung strikt), kappa = 1/2 (Gleichheit, Gauß), kappa = 1/4
% (Scheitel unter null — kann kein Zustand sein).
%
% ENTSCHEIDUNGEN
%   * die y-Achse ist so eng gefasst, dass der Höhenunterschied
%     kappa^2 = 1/4 vs 1/16 (also 0.19) noch gut 5 mm misst: die drei
%     Parabeln sind exakte Vertikalverschiebungen voneinander, weiter
%     lassen sie sich nicht trennen, und die Aussage steckt allein in
%     der Scheitelhöhe;
%   * schraffiert ist in (a) der Halbraum tr(rho M) < 0 und in (b) das
%     Gebiet Delta x Delta p < hbar/2 — zweimal dasselbe Verbot, einmal
%     als Operatoraussage, einmal als Aussage über Messstatistiken;
%   * Farbe kodiert das Produkt kappa und sonst nichts, dieselben drei
%     Farben in beiden Panels: die Parabel in (a) ist als Punkt in (b)
%     wiederzufinden.  Rot und gestrichelt ist der Fall, den es nicht
%     gibt;
%   * `axis on top' in (a): die Schraffur ist unbeschnitten (clip=false,
%     damit die Direktbeschriftung rechts stehen kann) und übermalt sonst
%     die Achslinie, die dann gestrichelt aussieht;
%   * grau ist ausschließlich Referenz: die Nulllinie, die Scheitel-
%     senkrechte lambda_*, und in (b) die Senkrechte Delta x = a, auf der
%     die drei Zustände aus (a) liegen;
%   * die beiden zusätzlichen Punkte auf der Hyperbel in (b) sind
%     dieselbe Gauß-Familie bei anderer Breite: quetscht man den Ort,
%     rutscht der Punkt die Hyperbel hinauf — das ist der Pfeil.  Die
%     Größenordnungsnotiz macht aus der willkürlichen Einheit a eine
%     Aussage über ein Elektron: auf einen Atomdurchmesser eingesperrt,
%     streut seine Geschwindigkeit um einige hundert km/s.
%
% Panel (a) die Parabelschar g(mu) = mu^2 + mu + kappa^2.
% Panel (b) die zulässigen Paare (Delta x, Delta p).

\begin{center}
\begin{tikzpicture}

% ─────────────────────────────────────────────────────────────
%  (a)  Die Positivitätsparabel
% ─────────────────────────────────────────────────────────────
\begin{axis}[
    name=parab,
    tpcplot,
    width=7.5cm, height=6.0cm,
    grid=none,
    xmin=-1.35, xmax=1.15,
    ymin=-0.40, ymax=1.78,
    xtick={-1,-0.5,0,0.5},
    xticklabels={$-1$,$-\frac12$,$0$,$\frac12$},
    ytick={0,0.5,1,1.5},
    xlabel={$\lambda\,(\Delta x)^{2}/\hbar$},
    ylabel={$\tr(\rho M)\,(\Delta x)^{2}/\hbar^{2}$},
    every axis x label/.style={at={(axis description cs:0.5,-0.15)},
                               anchor=north, font=\small},
    every axis y label/.style={at={(axis description cs:-0.14,0.5)},
                               rotate=90, anchor=south, font=\small},
    title={(a) \tpcifenglish{A parabola that may not dip}%
                            {Eine Parabel, die nicht eintauchen darf}},
    samples=181, clip=false, axis on top,
]

  % Verbotener Halbraum: M = C^dag C ist positiv, also tr(rho M) >= 0.
  \path[forbidden] (axis cs:-1.35,-0.40) rectangle (axis cs:1.15,0);
  \draw[levellinegrey] (axis cs:-1.35,0) -- (axis cs:1.15,0);

  % Scheitelsenkrechte mu = -1/2, also lambda_* = -hbar/2(Delta x)^2.
  \draw[classicallim] (axis cs:-0.5,-0.40) -- (axis cs:-0.5,1.10);

  % g(mu) = mu^2 + mu + kappa^2, kappa = Delta x Delta p / hbar.
  \addplot[wavefn, munsell,    domain=-1.35:0.35] {x^2 + x + 1};
  \addplot[wavefn, spacecadet, domain=-1.35:0.35] {x^2 + x + 0.25};
  \addplot[wavefn, birrenred, densely dashed, domain=-1.35:0.35]
    {x^2 + x + 0.0625};

  % Die Scheitel: das ist die ganze Information.
  \addplot[only marks, mark=*, mark size=2.0pt, munsell]
    coordinates {(-0.5,0.75)};
  \addplot[only marks, mark=*, mark size=2.0pt, spacecadet]
    coordinates {(-0.5,0)};
  \addplot[only marks, mark=*, mark size=2.0pt, birrenred]
    coordinates {(-0.5,-0.1875)};

  % Direktbeschriftung am rechten Kurvenende, in der Kurvenfarbe.
  \node[lbl, text=birrengray, font=\scriptsize, anchor=west]
    at (axis cs:0.42,1.66) {$\Delta x\,\Delta p$:};
  \node[lbl, text=munsell,    anchor=west] at (axis cs:0.42,1.4725)
    {$\hbar$};
  \node[lbl, text=spacecadet, anchor=west] at (axis cs:0.42,0.7225)
    {$\hbar/2$};
  \node[lbl, text=birrenred,  anchor=west] at (axis cs:0.42,0.5350)
    {$\hbar/4$};

  % Wo der Scheitel liegt und wie hoch er ist -- ein Block über der
  % Scheitelsenkrechten, im freien Feld über der munsell-Parabel.  Das
  % Minimieren über lambda, das die Vorlesung als Übung stellt, steht
  % damit als Bild da: der Scheitel sinkt mit dem Produkt.
  \node[lbl, text=birrengray, font=\scriptsize, anchor=south, align=center]
    at (axis cs:-0.5,1.12)
    {$\lambda_{*}(\Delta x)^{2}=-\hbar/2$\\[-1pt]
     \tpcifenglish{vertex height}{Scheitelhöhe}\\[-1pt]
     $(\Delta x\,\Delta p/\hbar)^{2}-\tfrac14$};

  % Warum der schraffierte Halbraum leer ist.
  \node[lbl, text=birrenred, font=\scriptsize, anchor=south east,
        fill=white, fill opacity=0.85, text opacity=1, inner sep=2pt]
    at (axis cs:1.13,-0.36)
    {$M=C^{\dag}C\ge 0$: \tpcifenglish{impossible}{unmöglich}};
\end{axis}

% ─────────────────────────────────────────────────────────────
%  (b)  Die zulässigen Paare (Delta x, Delta p)
% ─────────────────────────────────────────────────────────────
\begin{axis}[
    tpcplot,
    at={(parab.south east)}, anchor=south west, xshift=10mm,
    width=6.5cm, height=6.0cm,
    xmin=0, xmax=2.62, ymin=0, ymax=1.85,
    xtick={0,1,2}, ytick={0,0.5,1,1.5},
    grid=major,
    xlabel={$\Delta x/a$},
    ylabel={$\Delta p\,a/\hbar$},
    every axis x label/.style={at={(axis description cs:0.5,-0.15)},
                               anchor=north, font=\small},
    every axis y label/.style={at={(axis description cs:-0.125,0.5)},
                               rotate=90, anchor=south, font=\small},
    title={(b) \tpcifenglish{Which pairs survive}{Welche Paare übrig bleiben}},
    clip=false,
]
  % Verboten: Delta x Delta p < hbar/2.  Links von x = 1/(2 ymax) hat die
  % Hyperbel den Rahmen schon verlassen, dort wird der Streifen ganz
  % schraffiert.
  \path[forbidden] (axis cs:0,0) rectangle (axis cs:0.27027,1.85);
  \addplot[forbidden, draw=none, domain=0.27027:2.62, samples=200]
    {0.5/x} \closedcycle;

  % Die Schranke selbst.
  \addplot[spacecadet, line width=1.1pt, domain=0.27027:2.62, samples=200]
    {0.5/x};

  % Die Senkrechte Delta x = a trägt die drei Zustände aus (a).
  \draw[classicallim] (axis cs:1,0) -- (axis cs:1,1.02);

  \addplot[only marks, mark=*, mark size=2.2pt, munsell]
    coordinates {(1,1)};
  \addplot[only marks, mark=*, mark size=2.2pt, spacecadet]
    coordinates {(1,0.5)};
  % Ring thickness goes on the plot itself: nested inside mark options it
  % is typeset rather than applied, which shows up as nullfont warnings.
  \addplot[only marks, mark=o, mark size=2.6pt, birrenred, line width=1.0pt]
    coordinates {(1,0.25)};

  % Dieselbe Gauß-Familie bei anderer Breite: quetschen heißt die
  % Hyperbel entlangrutschen.
  \addplot[only marks, mark=*, mark size=2.2pt, spacecadet]
    coordinates {(0.5,1) (2,0.25)};
  \draw[-{Stealth[length=4pt]}, spacecadet, line width=0.7pt]
    (axis cs:1.05,0.60) .. controls (axis cs:0.76,0.74) ..
    (axis cs:0.60,0.96);

  \node[lbl, text=spacecadet, font=\scriptsize, anchor=south west]
    at (axis cs:0.58,1.08) {\tpcifenglish{squeeze $x$}{$x$ quetschen}};
  \node[lbl, anchor=south east, align=right] at (axis cs:2.58,0.40)
    {\tpcifenglish{Gaussian: equality}{Gauß: Gleichheit}\\[-1pt]
     $\Delta x\,\Delta p=\hbar/2$};

  % Das Verbot beschriftet den Streifen, in dem es gilt, statt die Mitte
  % der Tafel zu belegen.
  \node[lbl, text=birrengray, font=\scriptsize, rotate=90, anchor=center]
    at (axis cs:0.135,1.00) {\tpcifenglish{no state}{kein Zustand}};

  % Die Brücke zu Panel (a): gleiche Farben, gleiche drei Zustände.
  \node[lbl, text=birrengray, font=\scriptsize, anchor=west]
    at (axis cs:1.07,0.87) {\tpcifenglish{the three of (a)}{die drei aus (a)}};

  % Größenordnung, damit die Einheit a etwas bedeutet.
  \node[lbl, text=birrengray, font=\scriptsize, anchor=north east,
        align=right] at (axis cs:2.60,1.82)
    {$e^{-}$, $a=10^{-10}\,$m:\\[-1pt]
     $\Delta v\gtrsim 6\cdot10^{5}\,$m/s};
\end{axis}

\end{tikzpicture}
\end{center}
